\subsection{Transaction state and resource ownership}
A transaction-aware model must specify where a request resides while it is blocked. Merely delaying its completion is insufficient: the outstanding work may occupy an endpoint queue, a link credit, a memory slot, or none of these resources. We distinguish a producer-owned request that has not been admitted from an admitted record that consumes a modeled hardware queue entry. A failed admission leaves the former with its producer. Registering a wakeup callback does not itself allocate another hardware slot.

The implemented read path decomposes an operation into AR generation, fabric delivery, memory admission, overlapping access latency, bandwidth service, R generation, return transport, and destination consumption. Memory admission and response generation are distinct events. The read's useful-byte count belongs to returned data, not to the AR descriptor; adding both would double-count the logical payload. A network write instead carries an indivisible AW+W record and later generates a B response. A MoE combine is another write operation rather than the B acknowledgement for dispatch. This distinction matters when attributing response traffic to the application.

For a request $j$, let $a_j$ be arrival at the MPU, $m_j$ admission, $r_j$ latency readiness, $s_j$ first bandwidth service, and $e_j$ memory completion. The implemented service obeys
\begin{equation}
 a_j\leq m_j,\quad r_j=m_j+L_j,\quad s_j\geq r_j,\quad e_j\geq s_j.
\end{equation}
Admission waiting is $m_j-a_j$; latency waiting is $r_j-m_j$; bandwidth waiting is $s_j-r_j$. These intervals describe one request. Their sums over requests measure accumulated pressure, not an end-to-end application latency. Response transport and HBM service occur after or alongside other requests' intervals and require separate accounting.

This decomposition also exposes a limitation of the present model. At $e_j$, the memory slot is released even if the generated R record cannot immediately enter its output queue. The completed response remains pending in the adapter. Thus a memory-slot bound currently constrains memory service, not the entire request-to-response round trip. Applying a full network RTT to that slot bound would overstate its modeled restriction. Conversely, claiming a fully bounded memory subsystem would overlook pending response storage. A complete implementation must separately budget memory slots and completed-response bytes.

\subsection{Packetizer eligibility and ordering}
The packetizer operates on a per-output-port record queue. Its first queued record determines the packet's compatibility class. It scans queued records and selects eligible records of that class until the configured packet capacity is reached. Records of other classes can remain queued while compatible later records are selected. Selection therefore preserves an eligibility rule rather than a global FIFO over every transaction class. Evaluating a different selection rule requires an explicit policy change; it should not be hidden inside a purportedly equivalent implementation optimization.

The compatibility key includes source Die and port, destination chip and Die, destination proxy where applicable, VC, and record kind. This is stricter than equal destination chip alone. Two records sent toward different destination Dies or through different forwarding proxies cannot be merged merely because the same chip eventually owns the data. AW and W are represented as one eligible unit, preventing aggregation from separating a write address from its associated data. Each response retains its relation to the originating transaction even when several responses share one fabric packet.

After packet formation, the endpoint checks whether its device egress can admit the packet. If not, it leaves the selected records in their queue. On successful admission, it removes those records and reserves the next packet-formation opportunity using the modeled endpoint wire-slot duration, rounded to the configured whole-nanosecond opportunity schedule. The fabric can add further delay through arbitration, control traffic, and credit availability. Consequently, packetizer opportunities and actual link transmissions are different timestamps. The precision experiment does not imply that every endpoint scheduling boundary was converted into a fractional-nanosecond issue mechanism.

Best-effort packing is sensitive to simultaneous arrivals. A queue that has accumulated compatible records during backpressure can form a large packet without an explicit batching delay. Under light load, the same policy may send much smaller packets. The resulting packet-size distribution is an output of workload timing and contention, not a fixed user-provided average. Comparing policies using only a mean packet size would miss this feedback.

\subsection{A bounded completion path as a design extension}
Before evaluating aggressive credit allocation, we propose adding a finite completion-store resource at each MPU. A memory operation would need either a reserved response slot before service or available response bytes at completion. The design must choose and document one of these semantics. Early reservation prevents completion-stage overflow but can hold space throughout the access latency. Late reservation increases overlap but requires a place to retain an otherwise completed operation while waiting for storage.

Neither rule should silently increase physical capacity. The state machine must account for requests awaiting admission, in-service data, completed responses, and endpoint records without charging the same storage twice. If response storage is full, the corresponding backpressure should prevent an unbounded stream of newly admitted operations. This extension could use an explicit admission condition combining memory-slot availability and response-space reservation. It is not part of the reported experiments, so its impact on throughput or fairness remains unknown.

Correctness checks should verify conservation at every state transition: generated requests equal pending plus active plus completed requests; delivered useful bytes equal the sum of completed data fragments; and memory and response occupancy never exceed their independent limits. A final zero-pending check detects stuck operations but is insufficient on its own: a simulation can finish while having exceeded a buffer transiently. Peak and time-series assertions are therefore required as well.
